
\input texsis
\paper

\overfullrule=0pt
\def\frac#1#2{#1\over #2}
%
\def\bmatrix#1{\left[ \matrix{#1} \right]}
\def\be{{\beta}}
\def\toone{{t+1}}

%\def\frac#1#2{#1\over #2}
%\magnification=1200
\overfullrule=0pt

\line{Macro Dynamics \hfil NYUAD}
\line{Fall 2, 2018 \hfil Thomas Sargent}

\centerline{\bf Practice problems}
%\centerline{\bf Homework 3}


\medskip

\noindent{\bf Instructions:}  This document contains practice problems.  It is very likely that you will see one or more problems that very much
remind you of these problems on the in class final examination.
 


\medskip
\noindent{\it Exercise 3.1}
\medskip \noindent  Consider a consumer who wants to choose $\{c_t\}_{t=0}^T$ to maximize
$$\sum_{t=0}^T \beta^tc_t  \quad  , \;\;  0< \beta<1$$ subject to the intertemporal budget constraint
$$\sum_{t=0}^T R^{-t}[c_t-y_t]=0  \quad  , \;\;  R>1$$ where $c_t\geq0$ and $y_t>0$ for $t=0,\cdots,T$. Here $R>1$ is the gross interest rate $(1+r),\;\; r>0$. Assume $R$ is constant. Here $\{y_t\}_{t=0}^T $ is an exogenous sequence of ``labor income''.
\medskip
\noindent {\bf a.} Assume that $\beta R < 1$. Find the optimal path $\{c_t\}_{t=0}^T$.
\medskip
\noindent {\bf b.} Assume that $\beta R > 1$. Find the optimal path $\{c_t\}_{t=0}^T$.
\medskip
\noindent {\bf c.} Assume that $\beta R = 1$. Find the optimal path $\{c_t\}_{t=0}^T$.

\medskip
\noindent{\it Exercise 3.2}
\medskip
\noindent
 A representative consumer has preferences ordered by
 $$ \sum_{t=0}^\infty \beta^t \log(c_t), \;\;  0< \beta<1$$
 where $\beta = {\frac{1}{1+\rho}}, \rho>0$, and $c_t$ is the consumption per worker.
The technology is $$ y_t=f(k_t)=zk_t^\alpha, \;\; 0<\alpha<1, \;\; z>0 $$
where $y_t$ is the output per unit labor and $k_t$ is capital per unit labor
$$\eqalign{
y_t & =c_t+x_t+g_t \cr
  k_{t+1} & =(1-\delta)k_t+x_t, \;\;  0 < \delta<1 \cr
}$$
and $x_t$ is  gross investment per unit of labor and $g_t$ is  government expenditures per unit of labor.
Assume a competitive equilibrium with a price system $\{q_t,\eta_t,w_t\}_{t=0}^\infty$ and a government policy $\{g_t,T_t\}_{t=0}^\infty$, where 
$g_t$ is government consumption and $T_t$ is a lump sum tax at time $t$.

 Assume that the government finances its expenditures by levying lump sum taxes. There are no distorting taxes. Assume that at time 0, the economy begins
 with a capital per unit of labor $k_0$ that equals the steady state value appropriate for an economy in which $g_t$ had been zero forever.
\medskip
\noindent {\bf a.} Find a formula for the steady state capital stock when $g_t=0 \; \forall t$.
\medskip

\noindent {\bf b.} Compare the steady state capital labor ratio $\overline{k}$ in the competitive equilibrium with the capital labor ratio $\tilde{k}$ that maximizes steady state consumption per capita, i.e ,
 $\tilde{k}$ solves $$\tilde{c}=\max_k \left( f(k)-\delta k \right) $$
Is $\tilde{k}$ greater than or less that $\overline{k}$? If they differ, please tell why. Is $\tilde{c}$ greater or less than $\overline{c} = f(\overline{k})-\delta  \overline{k}$? Explain why.
\medskip

\noindent {\bf c.} Now assume that at time 0, $g_t$ suddenly jumps to the value $g={\frac{1}{2}}\overline{c}$ where
$\overline{c}$ is the value of consumption per capita in the initial steady state in which $g$ was zero forever. Starting from $k_0 = \overline{k}$ for the old $g=0$ steady state, find the time paths of $\{c_t,k_{t+1}\}_{t=0}^\infty$ associated with the new path $g_t=g>0 \; \forall t$ for government expenditures per capita.
Also show the time path for $\bar R_{t+1} \equiv (1-\delta) + f'(k_{t+1})$
Explain why the new paths are as they are.



\medskip
\noindent{\it Exercise 3.3}
\medskip
\noindent
 Assume 
 
 A representative consumer has preferences ordered by
 $$ \sum_{t=0}^\infty \beta^t \log(c_t), \;\;  0< \beta<1$$
 where $\beta = {\frac{1}{1+\rho}}, \rho>0$, and $c_t$ is the consumption per worker.
The technology is $$ y_t=f(k_t)=zk_t^\alpha, \;\; 0<\alpha<1, \;\; z>0 $$
where $y_t$ is the output per unit labor and $k_t$ is capital per unit labor
$$\eqalign{
y_t & =c_t+x_t+g_t \cr
  k_{t+1} & =(1-\delta)k_t+x_t, \;\;  0 < \delta<1 \cr
}$$
and $x_t$ is  gross investment per unit of labor and $g_t$ is  government expenditures per unit of labor.
Assume a competitive equilibrium with a price system $\{q_t,\eta_t,w_t\}_{t=0}^\infty$ and a government policy $\{g_t,T_t\}_{t=0}^\infty$, where
$g_t$ is government consumption and $T_t$ is a lump sum tax at time $t$.

 
Suppose that you observe the path for consumption per capita in figure 1. Say what you can about the likely behavior over time of $k_t$, $\bar R_t =[1 + (f'(k_t)-\delta)]$, $g_t$ and $T_t$. (You are free to make up any story that is consistent with the model.)

\medskip


\midfigure{fig_7_1}
\centerline{\epsfxsize=2.5truein\epsffile{fig_7_1.eps}}
\caption{Consumption per capita.}
%\infiglist{fig_7_1}
\endfigure




\medskip
\noindent{\it Exercise 3.4}
\medskip
\noindent Assume the same economic environment as in the previous  problem. Assume that someone has observed the  time path for $c_t$ in figure 2.
\medskip
\midfigure{fig_7_2}
\centerline{\epsfxsize=2.5truein\epsffile{fig_7_2.eps}}
\caption{Consumption per capita.}
%\infiglist{fig_7_2}
\endfigure
\medskip
\noindent {\bf a.} Describe a consistent set of assumptions about the fiscal policy that explains this time path for
$c_t$. In doing so, please distinguish carefully between changes in taxes and expenditures that are \underbar{foreseen} versus \underbar{unforeseen}.

\medskip
\noindent {\bf b.} Describe what is happening to $k_t,\bar R_t$ and $g_t$ over time. Make whatever assumptions you must to get a complete but consistent story - here ``consistent'' means ``consistent with the economic environment we have assumed''.


\midfigure{FN_2}
%\centerline{\epsfxsize=2.25truein\epsffile{FN_2_tom.eps}}
\centerline{\epsfxsize=2.25truein\epsffile{fig_11.4.eps}}
\caption{Capital stock as function of time in two economies with different values of $\gamma$.}
%\infiglist{FN_2}
\endfigure

\medskip
\noindent{\it Exercise 3.5}
\medskip
\noindent
 The structure of the  economy is identical to that described in the previous exercise.
Assume that $\{g_t, T_t\}_{t=0}^\infty $ are all constant sequences (their values don't change over time).
In this problem, we ask you to infer  differences
across two economies in which all aspects of the economy are identical {\it except} the parameter $\gamma$ in the utility function and maybe the level
of lump sum taxes $T_t$ imposed.  In both economies, $\gamma > 0$.  In one economy, $\gamma >0$ is high and in the other it is low.
Among other identical features, the two economies have identical government policies and identical initial capital stocks.\medskip
\medskip

\noindent{\bf a.}   Please look at figure 3. Please tell which outcome for $\{k_{t+1}\}_{t=0}^\infty$ describes the low $\gamma$ economy, and which describes the high $\gamma$ economy.
Please explain your reasoning.\medskip

\noindent{\bf b.}   Please look at figure 4. Please tell which outcome for $\{f'(k_t)\}_{t=0}^\infty$ describes the low $\gamma$ economy, and which describes the high $\gamma$ economy.
Please explain your reasoning.
\medskip

\noindent{\bf c.} Please plot time paths of consumption for the low $\gamma$ and the high $\gamma$ economies.
\medskip




\midfigure{FN_3}
%\centerline{\epsfxsize=2.25truein\epsffile{FN_3_tom.eps}}
\centerline{\epsfxsize=2.25truein\epsffile{fig_11_5.eps}}
\caption{Marginal product of capital as function of time in two economies with different values of $\gamma$.}
%\infiglist{FN_3}
\endfigure
\medskip

\medskip
\noindent{\it Exercise 3.6} \quad {\bf Trade and growth} \medskip
\medskip
\noindent
{\bf Part  I}
\medskip \noindent
Consider the problem of a planner in a small economy. When the economy is closed to international trade, the planner chooses $\{c_t,k_{t+1}\}_{t=0}^\infty$ to maximize
$$ \sum_{t=0}^\infty \beta^t u(c_t), $$
where $ 0< \beta<1, \beta = {\frac{1}{1+\rho}}, \rho>0$
subject to
$$c_t+k_{t+1}=f(k_t)+(1-\delta)k_t,  \quad  \delta \in (0,1)$$
\noindent where $u(c_t)={\frac{c_t^{1-\gamma}}{1-\gamma}}$, $\gamma>0$, $f(k_t)=zk_t^\alpha$, and $0<\alpha<1$.


\medskip
\noindent
Let $\overline{k}$ be the steady state value of $k_t$ under the optimal plan.
\medskip
\noindent{\bf a.} Find a formula for $\overline{k}$.
\medskip
\noindent{\bf b.} Assume that $k_0<\overline{k}$. Describe time paths for $\{c_t,k_{t+1}\}_{t=0}^\infty$ and $\bar R_{t+1}=(1-\delta) + f'(k_{t+1})$.
\medskip
\noindent{\bf c.} What is the steady state value of $\bar R_{t+1}$?\medskip

\noindent{\bf d.} Is $\bar R_{t+1}$ less or greater than its steady state value when $k_{t+1} < \overline{k}$?

\medskip\noindent
{\bf Part  II}
\medskip \noindent
Now assume that the economy is open to international trade in capital and financial assets. Assume that there is a fixed world gross rate of return $R=\beta^{-1}$ at which the planner can borrow or lend, what is often called a `small open economy' assumption.
The planner can use the proceeds of borrowing to purchase goods on the international market. These goods can be used to augment capital or to consume.\medskip \noindent
Let $\overline{k}_0$ be the level of initial capital $(\overline{k}_0<\overline{k})$ before the country opens up to trade just before time 0.  Let $\overline{k}_0$ be the same initial capital per capita $\overline{k}_0 < \overline{k}$ studied in parts {\bf a}-{\bf d}.\medskip \noindent
At time $t=-1$, after $\overline{k}_0$ was set, trade opens up. At time $t=-1$, the planner can issue IOU's or bonds in amount $B_{-1}$ and use the proceeds to purchase capital, thereby setting
$$k_0=\overline{k}_0+B_{-1} $$
where $B_{-1}$ is denominated in time -1 consumption goods. The bonds are one period in duration and bear the constant world gross interest rate $R=\beta^{-1}$.
For $t\geq 0$, the planner faces the constraints
$$ c_t+k_{t+1}+RB_{t-1}=f(k_t)+(1-\delta)k_t+B_t$$
Here $RB_{t-1}$ is the sum of interest and  principal on the bonds issued at $t-1$ and $B_t$ is the amount of one-period bonds issued at $t$. \medskip \noindent
The planning problem in the small open economy is now to choose $\{c_t,k_{t},B_t\}_{t=0}^\infty$ and $B_{-1}$, subject to $\overline{k}_0$ given. (Notice that $k_0$ is a choice variable and that $\overline{k}_0$ is an initial condition.)
Please solve the  planning problem in the small open economy and compare outcomes to those in the closed economy.
 Is  welfare $\sum_{t=0}^\infty \beta^t u(c_t)$ higher in the ``open'' or ``closed'' economy?







\noindent {\bf EXTRA CREDIT:}






\medskip
\noindent{\it Exercise 3.7}
\quad {\bf Trade and growth, version II} \medskip\noindent
{\bf Part  I}
\medskip \noindent
Consider the problem of the planner in a small economy. When the economy is \underbar{closed to trade}, the planner chooses $\{c_t,k_{t+1}\}_{t=0}^\infty$ to maximize
$$ \sum_{t=0}^\infty \beta^t u(c_t), \quad  0< \beta<1, \quad  \beta = {{\frac{1}{1+\rho}}}, \rho>0$$
subject to
$$c_t+k_t+1=f(k_t)+(1-\delta)k_t,  \delta \in (0,1)$$
\noindent where% $$\eqalign{u(c_t)={\frac{c_t^{1-\gamma}}{1-\gamma}} \quad  &, \;\; \gamma>0\cr
%f(k_t)=zk_t^\alpha \quad &,\;\;0<\alpha<1
%}$$
\medskip
\noindent
Let $\overline{k}$ be the steady state value of $k_t$ under the optimal plan.
\medskip
\noindent{\bf a.} Find a formula for $\overline{k}$.
\medskip
\noindent{\bf b.} Assume that $k_0 > \overline{k}$. Describe time paths for $\{c_t,k_{t+1}\}_{t=0}^\infty$ and $\bar R_{t+1}=(1-\delta) + f'(k_{t+1})$.
\medskip
\noindent{\bf c.} What is the steady state value of $\bar R_{t+1}$?\medskip

\noindent{\bf d.} Is $\bar R_{t+1}$ less or greater than its steady state value when $k_{t+1} > \overline{k}$?

\medskip\noindent
{\bf Part  II}
\medskip
\noindent{\bf e.} Now assume that the economy is open to international trade in capital and financial assets. Assume that there is a fixed world gross rate of return $R=\beta^{-1}$ at which the planner can borrow or \underbar{lend}.
In particular, the planner is free to use the following plan. The planner can sell all of its capital $\overline{k}_0$ and simply consume the interest payments.
\noindent
Let $\overline{k}_0$ be the level of initial capital $(\overline{k}_0>\overline{k})$ just before the country opens up to trade just before time 0.  Let $\overline{k}_0$ be the initial capital per capita $\overline{k}_0 >\overline{k}$ studied in parts {\bf a}--{\bf d}.
At time $t=-1$, after $\overline{k}_0$ was set, trade opens up. At time $t=-1$, the planner \underbar{sells} $\overline{k}_0$ in exchange for IOU's or bonds from the rest of the world in the amount $A_{-1}=\overline{k}_0$.
\medskip \noindent
The bonds $A_{-1}$ are one-period in duration and bear the constant world gross interest rate $R=\beta^{-1}$.
After the sale of $\overline{k}_0=A_{-1}$, the planner has \underbar{zero} capital and so shuts down the technology. Instead, the planner uses the asset market to smooth consumption. The planner chooses $\{A_{t+1},c_t\}$ to maximize
$$ \sum_{t=0}^\infty \beta^t u(c_t),  \quad   0< \beta<1,\quad   \beta = {\frac{1}{1+\rho}},  \quad \rho>0$$
subject to
$c_t+\beta A_{t+1}=A_t,  A_0 = \overline{k}_0 .$
Please find the optimal path for consumption $\{c_t\}$.
%                   Form the Lagrangian
%$$L = \sum_{t=0}^\infty \beta^t \{u(c_t) + \lambda_t[\beta^{-1}(A_t-c_t)-A_{t+1}]\}$$
%
% $$\eqalign{ FONC &\colon \cr
%c_t &\colon\; u'(c_t)-\lambda_t\beta^{-1}=0 \cr
%A_t &\colon\; -\beta^-1\lambda_{t-1} + \beta^{-1}\lambda_t=0 \cr
%&\lim_{T\to\infty}\beta^{T}\lambda_T A_{T+1}=0 \cr
%\Rightarrow\;\; &\lambda_t  =\lambda_0 \;\; \forall t \cr
%&u'(c_t) =u'(c_0)\;\; \forall t \cr
%&c_t =c_0\;\; \forall t \cr
%$The$&$n rollover assets and just consume interest:$ \cr
%&A_{t+1} =\beta^{-1}[A_t -c_t] \Rightarrow \cr
%&\bar A =\beta^{-1}[\bar A -c_0] \cr
%&(1-\beta^{-1})\bar A = -c_0 \cr
%& \quad $or$ \cr
%& c_0  =(\beta^{-1}-1)\bar A \cr
%& c_0  =(1+\rho-1)\bar A \cr
%& c_0 =\rho \bar A \quad, \;\; \bar A= A_0=\overline{k}_0 \cr\
%}$$
%\medskip
%\noindent
%Thus, the optimal path for consumption is c_t = \rho \bar A=\rho \overline{k}_0
\medskip
\noindent{\bf f.} Compare the path of $(c_t,k_{t+1},\bar R_t)$ that you computed in parts {\bf a}-{\bf d} with ``no trade'' with the part {\bf e} path ``with trade'' in which the government ``shuts down the home technology'' and lives entirely from the returns on foreign assets. Can you say which path the representative consumer will prefer?
%\medskip \underbar{Hint:} I think it could go either way for reasons that might be easier to appreciate after you look at part g.
\medskip
\noindent{\bf g.} Now return to the economy in part {\bf e} with $\overline{k}_0>\bar k$ from part {\bf d}.  Assume that the planner is free to borrow or lend capital at the fixed gross interest of $\beta^{-1}=1+\rho$, as before. But now assume the planner chooses the \underbar{optimal} amount of $\overline{k}_0$ to sell off and so does not necessarily sell off the entire $\overline{k}_0$ and possibly continues to operate the technology.

\medskip
\indent {\bf i)} Find the solution of the planning problem.\medskip
%\underbar{Hint:} You might want to review your solution to problem 2 of the practice problem 4.
%\medskip
\indent {\bf ii)} Explain why it is optimal not to shut down the technology. {\bf Hint}: Starting from having shut the technology down, think of putting a small amount $\epsilon$ of capital into the technology-this earns $z\epsilon^\alpha$ and costs $\rho \epsilon$ in terms of foregone interest. Because $0<\alpha<1$, $\rho\epsilon<z\epsilon^\alpha$ for small $\epsilon$. Thus, the technology is \underbar{very productive} for small $\epsilon$, so it is efficient to use it.

%\vfil\eject



\bye






\bye
